\centerline{\bf \Huge Муниципальная комба}

\problem{1}
$\bigl[\textbf{Новосибирск, 24/25, 9.2.}\bigr]$
В классе учатся 24 школьника, каждый из которых либо лгун (всегда лжёт), либо правдун (всегда говорит правду). На вопрос учителя: «Верно ли, что в вашем классе поровну лгунов и правдунов?» некоторые ответили «Да», а остальные - «Нет». Как учителю, зная ответ каждого ученика, определить, кто из них лгун, а кто - правдун?


\problem{2}
$\bigl[\textbf{Курская область, 17/18, 9.4.}\bigr]$
В компании из $2 n+1$ человек для любых $n$ человек найдется отличный от них человек, знакомый с каждым из них. Докажите, что в этой компании есть человек, знающий всех.

\problem{3}
$\bigl[\textbf{Москва, 20/21, 10.6.}\bigr]$
На острове живут два племени: рыцари и лжецы. Рыцари всегда говорят правду, а лжецы всегда лгут. Однажды 80 человек сели за круглый стол, и каждый из них заявил: «Среди 11 человек, сидящих следом за мной по часовой стрелке, есть хотя бы 9 лжецов». Сколько рыцарей сидит за круглым столом? Укажите все возможные варианты.

\problem{4}
$\bigl[\textbf{Республика Марий Эл, 20/21, 10.5.}\bigr]$
На доске $10 \times 10$ расставлены 10 фишек так, что в каждом горизонтальном и вертикальном ряду стоит по одной фишке. Можно ли оставшуюся часть доски замостить прямоугольниками $1 \times 2$ (по клеткам)?


\problem{5}
$\bigl[\textbf{Москва, 20/21, 11.5.}\bigr]$
В шахматном турнире соревнуются друг с другом команда школьников и команда студентов, в каждой из которых по 15 человек. В течение турнира каждый школьник должен сыграть с каждым студентом ровно один раз, причём каждый день каждый человек должен играть не более одного раза. В разные дни могло проводиться разное количество игр.

В некоторый момент турнира организатор заметил, что может составить расписание на следующий день из 15 игр ровно 1 способом, а из 1 игры $-N$ способами (порядок игр в расписании не важен, важно лишь кто с кем играет). Найдите наибольшее возможное значение $N$.

\newpage

\hspace{3mm}

\problem{6}
$\bigl[\textbf{Свердловская область, 18/19, 11.6.}\bigr]$
В однокруговом шахматном турнире (каждый шахматист играет с каждым одну партию) участвовало 20 шахматистов, причём 6 из них - из России. Известно, что набрав очков больше, чем кто-либо, первое место занял россиянин Владимир. Второе место занял Левон из Армении, также опередив по очкам каждого из остальных 18 шахматистов. Какое наибольшее суммарное количество очков могли набрать российские шахматисты? (В шахматах за победу в партии даётся одно очко, за ничью - пол-очка, за поражение очков не дают.)